% [原文 p.169]
\begin{proof}[证明（续）]
这样便知 $\Pic Y_n\to\Pic Z_n(0)$ 是满射，因此只需再证 $\Pic Z_n(0)\to\Pic\Gamma$ 是满射。

设 $S$ 是任意非奇异曲线，$\Gamma\subset S$ 是曲线的 $A_n$、$D_n$ 或 $E_n$ 型图。假设 $S$ 中有曲线 $C$，它与 $\Gamma$ 只在一条端曲线 $E$ 处相交，且 $(E\cdot C)=1$。在 $D$、$E$ 型的情形，再假设有曲线 $C'$，它与 $\Gamma$ 只在另一条端曲线 $E'$（$E'\neq E$）处相交，且 $(E'\cdot C')=1$。此时容易看出 $\Pic S\to\Pic\Gamma$ 是满射。（有不需要 $C'$ 的情形。此外，$E_8$ 的情形连 $C$ 都不需要。）

把这个结论用于 (10) 的证明：取 $Z_n(0)\to Z$ 的 $(-1)$-曲线作为 $C$；取 $r_n(0)^*((\text{一般点})\times\PP^1)$（$D$ 型）或 $r_n(0)^*(\text{一般直线})$（$E$ 型）作为 $C'$ 即可。这就证明了 (10)。
\end{proof}

\begin{proof}[4.41（定理 4.28 的证明）]
只须验证命题 4.31 的条件成立。$\bar p_n\colon\bar Y_n\to S$ 已经构造出来。由定理 4.39(7)–(9)，存在带有 $A_n$、$D_n$ 或 $E_n$ 型奇点的纤维。设 $x\in\bar Y_n$ 是对应于该奇点的点，则可知 $x$ 处的芽满足命题 4.31 的条件。事实上，命题 4.31(1) 由构造成立，命题 4.31(2) 来自例 4.30，命题 4.31(3) 来自定理 4.39(7)–(9)，而命题 4.31(4) 由定理 4.39(6) 导出。
\end{proof}

\begin{remark}{4.42}
由前述 Du Val 奇点形变的具体描述，可以简单地刻画邻近纤维中可能出现的奇点。例如，由此容易构造出从方程式的观点看并不明显的退化 $D_n\leftarrow A_{n-1}$、$E_n\leftarrow D_{n-1}$、$E_n\leftarrow A_{n-1}$。
\end{remark}

这里的论证看不出下面这一点，但定理 4.28 的构造可以整体化。后文并不需要，但为完整起见，在此陈述其结果。

\begin{theorem}{4.43}
（[Bri71, Art74]）设 $f\colon X\to S$ 是曲面的平坦族，使每个纤维 $X_s$ 都只含 Du Val 奇点。
% [原文 p.170]
则存在有限满射 $g\colon S'\to S$ 与同时奇点消解 $g^*f\colon X\times_S S'\to S'$：
\begin{tikzcd}
\bar X' \arrow[r,"p"] \arrow[d,"\bar f'"'] & X\times_S S' \arrow[r,"f^*g"] \arrow[d,"g^*f"] & X \arrow[d,"f"] \\
S' \arrow[r,equals] & S' \arrow[r,"g"'] & S.
\end{tikzcd}
\end{theorem}
% [Tutor Note: 有限基变换 $g$ 不可省：Du Val 奇点族的整体同时消解一般不存在，须先作有限基变换 $S'\to S$ 才能消解；这正是 §4.3 构造所反映的现象。]

\section{二维椭圆奇点}

\takeaway{本节以极小消解 $f\colon Y\to X$ 的例外循环 $Z$ 为桥，把 $Y$ 上的上同调与线丛截面问题化为 $Z$ 及其零维子概型上的问题，并由此得到椭圆奇点的标准方程（定理 4.57）。}

本节的目的是介绍二维椭圆奇点的构造理论。这些结果在定理 5.35 中研究三维标准奇点时是必要的。主要定理 4.57 属于 [Lau77, Rei76]。

\begin{notation}{4.44}
本节中，$(0\in X)$ 表示 $0$ 是唯一奇点的正规曲面芽。代数情形与解析情形没有区别。通常把 $X$（选取代表）看作簇，但选取 $X$ 时要使其在代数情形下为仿射的、在解析情形下为 Stein 的。

$f\colon Y\to X$ 为极小奇点消解，$E=\sum E_i$ 为例外曲线。与记号 4.1 同样，可以写成
\[ K_Y\equiv\sum a_iE_i\qquad (a_i\le 0). \]
记 $Z:=\sum[-a_i]E_i$、$\Delta_Y:=\sum\{-a_i\}E_i$，则 $K_Y+\Delta_Y\equiv -Z$ 成立。$X$ 是对数终端（lt）的当且仅当 $Z=0$。即使 $X$ 不是 lt，也可能出现 $\Supp Z\neq E$。若 $K_X$ 是 Cartier 的且 $X$ 不是 Du Val，则由推论 4.3 有 $\Supp Z=E$。\hfill$\square$
\end{notation}
% [Tutor Note: $Z=\sum[-a_i]E_i$ 与 $\Delta_Y=\sum\{-a_i\}E_i$ 是把 $K_Y$ 拆成整部与分数部；$K_Y+\Delta_Y\equiv -Z$ 是本节一切计算的出发点，$Z$ 还是数值 $f$-平凡的循环。]

首先说明，$Y$ 的若干上同调性质的研究在许多情形下可归结为 $Z$ 的研究。

\begin{propositionp}{4.45}
记号同上。设 $L$ 是 $Y$ 上 $f$-nef 的线丛。则下列成立。
\begin{enumerate}[label=(\arabic*),leftmargin=2.4em,itemsep=1pt,topsep=2pt]
\item $H^0(Y,L)\to H^0(Z,L|_Z)$ 是满射。
\item $H^1(Y,L)\cong H^1(Z,L|_Z)$。
\item $L\cong\cO_Y$ 当且仅当 $L\equiv_f 0$ 且 $L|_Z\cong\cO_Z$。
\item $f_*\omega_Y(Z)=\omega_X$。
% [原文 p.171]
\item $\omega_X/f_*\omega_Y\cong H^0(Z,\omega_Z)$。
\item $H^1(Y,\cO_Y)$ 与 $\omega_X/f_*\omega_Y$ 互相对偶。
\end{enumerate}
\end{propositionp}

\begin{proof}
考虑正合列
\[ 0\to L(-Z)\to L\to L|_Z\to 0, \]
并对其应用 $f_*$。由 $L(-Z)\equiv K_Y+\Delta_Y+L$ 与命题 2.71 得 $R^1f_*L(-Z)=0$。由此得到 (1) 与 (2)。

若 $Z=0$，则 $X$ 是 lt 的，(3) 由引理 4.13 推出。下面设 $L\equiv_f 0$、$Z\neq 0$ 且 $L|_Z\cong\cO_Z$。由 (1)，$L|_Z\cong\cO_Z$ 的常数截面 1 可延拓为 $L$ 的截面 $\sigma$。由 $L\equiv_f 0$，$\sigma$ 在 $E$ 附近生成 $L$。于是 $L\cong\cO_Y$。

$\omega_X$ 的任意截面给出 $\omega_Y(Z)$ 的有理截面，它至多沿 $E$ 有极点。因此 (4) 可在引理 4.46 中取 $L=K_Y+Z$、$\Delta=\Delta_Y$、$M=0$ 得到。由此可知 (5) 成立，这从下面的正合列与推论 2.68 可以看出：
\[ 0\to f_*\omega_Y\to f_*\omega_Y(Z)\to H^0(Z,\omega_Z)\to R^1f_*\omega_Y=0. \]
最后，(6) 是 (2) 与 (5) 的推论。
\end{proof}
% [Tutor Note: (6) 是消解上 Serre 对偶的局部形式；(0 是椭圆奇点) 时 $H^1(Y,\cO_Y)\cong\C$，维数 1 正是“椭圆”一词的来源。]

\begin{lemmap}{4.46}
记号同上。设 $L$ 是 $Y$ 上的线丛，且 $L\equiv -M-\Delta$。其中 $M$ 是 $f$-nef 的，$[\Delta]=0$，$\Supp\Delta\subset E$。则 $L|_{Y\setminus E}$ 的任意截面 $\sigma$ 都可延拓为 $L$ 的截面。
\end{lemmap}

\begin{proof}
$\sigma$ 是 $L$ 的有理截面，故可写成 $L\cong\cO_Y(B+C)$。其中 $B$ 是满足 $\Supp B\subset E$ 的除子，$C$ 是每个不可约分支都在 $X$ 上有限的有效除子。需证 $B$ 是有效的。由 $-B-\Delta\equiv M+C$，且 $M+C$ 是 $f$-nef 的。故由引理 3.41，$-\Delta\le B$，特别地 $B$ 有效。
\end{proof}

利用循环 $Z$，可以给出 $K_X$ 成为 Cartier 的数值判别。

\begin{propositionp}{4.47}
记号同 4.44。则下列条件等价。
\begin{enumerate}[label=(\arabic*),leftmargin=2.4em,itemsep=1pt,topsep=2pt]
\item $K_X$ 是 Cartier 的。
\item $\Delta_Y=0$，且对每个有效除子 $Z'\subseteq Z$ 都有 $h^1(\cO_{Z'})<h^1(\cO_Z)$。
\end{enumerate}
\end{propositionp}

\begin{proof}
若 $\Delta_Y=0$ 且 $Z=0$，则已知 $X$ 是 Du Val 的，
% [原文 p.172]
故可设 $Z\neq 0$。又若 $K_X$ 是 Cartier 的，则 $\Delta_Y=0$，故可设 $\Delta_Y=0$。对每个 $Z'\subseteq Z$ 考虑正合列
\[ 0\to\omega_Y\to\omega_Y(Z')\to\omega_{Z'}\to 0, \]
由 $H^1(\omega_Y)=0$（推论 2.68）可导出
\begin{equation}
H^0(Y,\omega_Y(Z'))/H^0(Y,\omega_Y)\cong H^0(Z',\omega_{Z'}). \tag{4.1}
\end{equation}
因此
\begin{align*}
h^1(\cO_{Z'})<h^1(\cO_Z) &\iff h^0(\omega_{Z'})<h^0(\omega_Z) \\
&\iff H^0(\omega_Y(Z'))\subsetneq H^0(\omega_Y(Z)).
\end{align*}
其中第一个等价由对偶得到，第二个由同构 (4.1) 得到。最后的断言对所有 $Z'\subseteq Z$ 成立，等价于 $\omega_Y(Z)$ 在 $Z$ 的所有泛点处由全局截面生成。由于 $\omega_Y(Z)$ 数值 $f$-平凡，这等价于 $\omega_Y(Z)\cong\cO_Y$。

若 $K_X$ 是 Cartier 的，则 $\omega_Y(Z)\cong f^*\cO_X(K_X)\cong\cO_Y$。反之设 $\omega_Y(Z)\cong\cO_Y$。此时 $\omega_X=f_*\omega_Y(Z)\cong f_*\cO_Y=\cO_X$，故 $K_X$ 是 Cartier 的。
\end{proof}

\begin{definition}{4.48}
记号同 4.44。当 $K_X$ 是 Cartier 的且 $R^1f_*\cO_Y\cong\C$ 成立时，称 $(0\in X)$ 为椭圆奇点（elliptic singularity）。由命题 4.45，后一条件等价于 $f_*\omega_Y=m_{0,X}\omega_X$。
\end{definition}

\begin{lemmap}{4.49}
假设 $(0\in X)$ 是椭圆的，$Z$ 如记号 4.44。若 $Z$ 是不可约且约化的曲线，则 $p_a(Z)=1$。其他情形下，$E$ 的每个不可约分支 $E_i$ 都是满足 $(E_i\cdot(-Z+E_i))=-2$ 的非奇异有理曲线。
\end{lemmap}

\begin{proof}
由命题 4.45(2)，$H^1(Y,\cO_Y)\cong H^1(Z,\cO_Z)$，故若 $Z$ 不可约且约化则证毕。否则 $E_i\subsetneq Z$，由命题 4.47 得 $h^1(\cO_{E_i})<h^1(\cO_Z)=1$。故 $h^1(\cO_{E_i})=0$，$E_i\cong\PP^1$。最后断言由 $-Z\equiv K_Y$ 取伴随（adjunction）得到。
\end{proof}
% [Tutor Note: $Z$ 不可约且约化时 $p_a(Z)=1$ 是“椭圆”几何含义所在；一般情形则由 $(-2)$-曲线链构成，与 Du Val 情形对照可见 $Z$ 刻画了奇点的“非 Du Val 部分”。]

% [原文 p.173]
至此，我们已经把关于 $Y$ 的若干问题转化成了关于 $Z$ 的问题。下面为了研究 $Y$ 上线丛的截面，先把问题归结到 $Z$，再进一步归结到 $Z$ 的零维子概型。最终目标是理解形如 $\bigoplus_{m\ge 0}H^0(Y,L^m)$ 的环。

\begin{lemmap}{4.50}
设 $V$ 是（不必约化的）固有曲线且满足 $H^1(\cO_V)=0$，$L$ 是 $V$ 上的 nef 线丛。则下列成立。
\begin{enumerate}[label=(\arabic*),leftmargin=2.4em,itemsep=1pt,topsep=2pt]
\item $L$ 由全局截面生成。
\item $H^1(V,L)=0$。
\end{enumerate}
\end{lemmap}

\begin{proof}
设 $V_i$ 是 $\operatorname{red}(V)$ 的不可约分支，$m_i=\deg(L|_{V_i})$。选取一般点 $P_i\in V_i$ 与 Cartier 除子 $D_i\subset V$，使 $D_i\cap V_i=P_i$。此时 $L':=\cO_V(\sum m_iD_i)$ 是 $V$ 上的线丛，且 $L$ 与 $L'$ 在每个 $V_i$ 上次数相同。故由引理 4.13 有 $L\cong L'$。这说明 $L$ 在 $P_i$ 之外由全局截面生成；由于 $P_i$ 可以移动，得到 (1)。

(2) 由对满射 $H^0(V,L)\otimes\cO_V\to L$ 取 $H^1$ 得到。
\end{proof}

\begin{propositionp}{4.51}
$(0\in X)$ 是椭圆的，$Z$ 如记号 4.44。设 $L$ 是 $Z$ 上的 nef 线丛且 $(L\cdot Z)>0$。则下列成立。
\begin{enumerate}[label=(\arabic*),leftmargin=2.4em,itemsep=1pt,topsep=2pt]
\item $H^1(Z,L)=0$。
\item 可以通过适当选取截面 $s\in H^0(Z,L)$，使得 $(s=0)$ 是与 $\operatorname{red}(Z)$ 的奇点集合不相交的零维子概型，且 $s|_{\operatorname{red}(Z)}$ 没有重零点。
\item 设 $C\subset Z$ 是 $Z$ 的不可约分支且 $(L\cdot C)>0$，令 $Z'=Z-C$。则 $H^0(Z,L)\to H^0(Z',L|_{Z'})$ 是满射。
\end{enumerate}
\end{propositionp}

\begin{proof}
若 $Z$ 不可约且约化，则由引理 4.49 得到命题 4.51。否则按 (3) 选取 $C$、$Z'$。于是由命题 4.47 有 $H^1(\cO_{Z'})=0$，由引理 4.50 知 $L|_{Z'}$ 由全局截面生成，且 $H^1(Z',L|_{Z'})=0$。考虑正合列
\[ 0\to L(-Z')|_C\to L\to L|_{Z'}\to 0. \]
由引理 4.49，$C\cong\PP^1$ 且 $\deg L(-Z')|_C=(L\cdot C)-2\ge -1$。故 $H^1(C,L(-Z')|_C)=0$，由相应的上同调列
% [原文 p.174]
得到 $H^1(Z,L)=0$。因此 (1) 成立。并且 $H^0(Z,L)\to H^0(Z',L|_{Z'})$ 是满射，得到 (3)。注意 $\operatorname{red}(Z)$ 的奇点集合含于 $\operatorname{red}(Z')$ 中。由此，$H^0(Z',L|_{Z'})$ 的一般元可延拓为 $s\in H^0(Z,L)$，并可使得 $(s=0)\cap Z'$ 是与 $\operatorname{red}(Z)$ 的奇点集合不相交的零维子概型。特别地，$(s=0)\not\supset C$，故由 Bertini 定理得到 (2)。
\end{proof}

\begin{notation}{4.52}
$(0\in X)$ 是二维椭圆奇点，$L$ 是 $Z$ 上满足 $(L\cdot Z)>0$ 的 nef 线丛，截面 $s\in H^0(Z,L)$ 选取为使 $\{s=0\}$ 是与 $\operatorname{red}(Z)$ 的奇点集合不相交的零维子概型、且 $s|_{\operatorname{red}(Z)}$ 没有重零点。

记 $V:=(s=0)$。此时 $A:=\cO_V$ 是半局部环，其根为 $m$，有分解 $A=\bigoplus_{i=1}^r A_i$、$m=\bigoplus_{i=1}^r m_i$。其中 $(A_i,m_i)$ 是局部 Artin $\C$-代数。记 $V_i=\Spec A_i$，则 $V_i\subset Z$ 是 Cartier 除子，$V=\sum V_i$ 也是到连通分支的分解。$\operatorname{socle}(m):=\{a\in m\mid ma=0\}$ 称为 $m$ 的基座（socle）。注意若 $A_i=\C[t]/(t^a)$，则 $\operatorname{socle}((t))=(t^{a-1})$。

记
\[ W_L:=\operatorname{im}[H^0(Z,L)\to A\otimes L]\subset A\otimes L. \]
$W_L$ 是 $A\otimes L$ 的 $\C$-向量子空间，但一般不是 $A$-子模。\hfill$\square$
\end{notation}

研究像 $W_L$ 是理解 $Y$ 上线丛的全局截面的重要线索。

\begin{lemmap}{4.53}
记号与假设同 4.52。则下列成立。
\begin{enumerate}[label=(\arabic*),leftmargin=2.4em,itemsep=1pt,topsep=2pt]
\item 若 $A\neq 0$，则 $\operatorname{codim}_{A\otimes L}W_L=1$。
\item 对每个 $j$，射影 $W_L\to A\otimes L\to(A/A_j)\otimes L$ 是满射。
\item 若 $m\neq 0$，则射影 $W_L\to A\otimes L\to(A/\operatorname{socle}(m))\otimes L$ 是满射。
\item 若 $\dim_{\C}A\ge 2$，则 $W_L$ 作为 $A$-模生成 $A\otimes L$。
\end{enumerate}
\end{lemmap}

\begin{proof}
由 $L=\cO_Z(V)$ 与 $\cO_V(V)=A\otimes L$ 得到正合列
\[ 0\to\cO_Z\to L\to A\otimes L\to 0, \]
% [原文 p.175]
由此列得到正合列 $H^0(Z,L)\to A\otimes L\to H^1(\cO_Z)\to H^1(L)$。其中 $H^1(\cO_Z)\cong\C$，且由命题 4.51(1) 也有 $H^1(L)=0$。由此得到 (1)。

对每个 $j$，由 $L=\cO_Z(V)$ 与 $\cO_{V-V_j}(V)=(A/A_j)\otimes L$ 得到正合列
\[ 0\to\cO_Z(V_j)\to L\to(A/A_j)\otimes L\to 0. \]
由命题 4.51(1) 有 $H^1(Z,\cO_Z(V_j))=0$，故得到 (2)。

其次设 $m\neq 0$，即 $\operatorname{socle}(m)\neq A$。设 $C\subset Z$ 是不可约曲线，且 $(s=0)$ 在 $C$ 上有非约化点。此时 $(L\cdot C)>0$，且由记号 4.52，$Z$ 沿 $C$ 不是既约的。记 $Z'=Z-C$。设 $s'$ 是 $s$ 到 $L|_{Z'}$ 的限制，$D':=(s'=0)\subset Z'$。得到正合列
\[ 0\to\cO_{Z'}\xrightarrow{s'}L|_{Z'}\to L|_{D'}\to 0, \]
由命题 4.47 得到满射 $H^0(Z',L|_{Z'})\to H^0(D',L|_{D'})$。与命题 4.51(3) 的满射合成，$H^0(Z,L)\to H^0(D',L|_{D'})$ 成为满射。$(A/\operatorname{socle}(m))\otimes L$ 是 $H^0(D',L|_{D'})$ 的商向量空间，故得到 (3)。

若 (4) 不成立，则 $W_L$ 的所有元都在 $\Spec A$ 的某点处消失。但由 (2) 与 (3) 这是不可能的。
\end{proof}
% [Tutor Note: $W_L$ 把 $Y$ 上的全局截面压到零维 Artin 环 $A$ 上，$A$ 的每个连通分支都是 $\C[t]/(t^a)$，于是截面的障碍完全变成线性代数；(4) 说明只要 $\dim_{\C}A\ge 2$，$W_L$ 就足够大。]

从技术层面说，下面是本节的主要结果。

\begin{propositionp}{4.54}
（[Lau77, Rei76]）$(0\in X)$ 是二维椭圆奇点，$Z$ 如记号 4.44。设 $L$ 是 $Z$ 上的 nef 线丛，$k=(L\cdot Z)$。则下列成立。
\begin{enumerate}[label=(\arabic*),leftmargin=2.4em,itemsep=1pt,topsep=2pt]
\item 若 $k\ge 2$，则 $L$ 由全局截面生成。
\item 若 $k\ge 3$，则 $\bigoplus_{n\ge 0}H^0(Z,L^{\otimes n})$ 由次数 1 的元生成。精确地说，
\[ \bigoplus_{n\ge 0}H^0(Z,L^{\otimes n})\cong\C[x_1,\dots,x_k]/I \]
成立。其中 $\deg x_i=1$，$I$ 由次数 2 与次数 3 的元生成。
% [原文 p.176]
\item 若 $k=2$，则 $\bigoplus_{n\ge 0}H^0(Z,L^{\otimes n})$ 由次数 1 与次数 2 的元生成。精确地说，
\[ \bigoplus_{n\ge 0}H^0(Z,L^{\otimes n})\cong\C[x,y,z]/(z^2+q_4(x,y)) \]
成立。其中 $\deg(x,y,z)=(1,1,2)$，$q_4$ 是次数 4 的齐次多项式或 0。
\item 若 $k=1$，则 $\bigoplus_{n\ge 0}H^0(Z,L^{\otimes n})$ 由次数 1、2、3 的元生成。精确地说，
\[ \bigoplus_{n\ge 0}H^0(Z,L^{\otimes n})\cong\C[x,y,z]/(z^2+y^3+ayx^4+bx^6) \]
成立。其中 $\deg(x,y,z)=(1,2,3)$，$a,b\in\C$。
\end{enumerate}
\end{propositionp}

\begin{proof}
记 $R_Z(n):=H^0(Z,L^{\otimes n})$，$R_Z:=\bigoplus_{n\ge 0}R_Z(n)$ 是对应的分次 $\C$-代数。取 $s\in H^0(Z,L)$ 为如记号 4.52 中的一般元。取 $V=(s=0)$、$A=\cO_V$、$(A,m)=\bigoplus_{i=1}^r(A_i,m_i)$ 如记号 4.52。此时由 $A$ 的定义，$\dim_{\C}A=k$。记 $R_V(n):=H^0(V,L^{\otimes n})$，$R_V:=\bigoplus_{n\ge 0}R_V(n)$。注意：若 $T$ 是生成 $L|_V$ 的 $L|_V$ 的任意截面，则 $R_V(n)=AT^n$ 且 $R_V=A[T]$ 成立。

注意正合列
\[ 0\to R_Z(n-1)\xrightarrow{s}R_Z(n)\to R_V(n)\to H^1(Z,L^{n-1})\to H^1(Z,L^n). \]
由命题 4.51(1)，$R_Z/sR_Z$ 是 $A[T]$ 的分次 $\C$-子代数，且满足
\[
R_Z/sR_Z(n)=\begin{cases}
AT^n & n\ge 2\ \text{时},\\
W_L & n=1\ \text{时},\\
\C & n=0\ \text{时}.
\end{cases}
\]
因此若 $k=1$，则 $A=\C$，且
\[ R_Z/sR_Z\cong\C[T^2,T^3]\cong\C[y,z]/(y^3-z^2) \]
成立，其中 $\deg y=2$、$\deg z=3$。于是 (4) 由引理 4.55 经显然的坐标变换得到。

% [原文 p.177]
假设 $k\ge 2$。由引理 4.53(4)，$W_L$ 生成 $A\otimes L$，故 (1) 得证。此后可以假定上面用到的 $T$ 是 $L$ 的某个全局截面 $t$ 的限制。于是 $t\in W_L$，且对 $n\ge 2$ 有
\[ t(R_Z/sR_Z(n))=(R_Z/sR_Z)(n+1). \]

因此若 $k=2$，则容易看出 $R_Z/(s,t)=\C\oplus\C z$，其中 $\deg z=2$。于是 $R_Z/(s,t)\cong\C[z]/(z^2)$，而 (3) 由引理 4.55 得到。

假设 $k\ge 3$。假设 $R_Z$ 不由 $R_Z(1)$ 生成，导出矛盾。由引理 4.55，$R_Z/sR_Z$ 也不由 $R_Z/sR_Z(1)$ 生成。进一步注意
\[ t(R_Z/sR_Z(n))=(R_Z/sR_Z)(n+1)\qquad(n\ge 2), \]
这意味着
\[ t(R_Z/sR_Z(1))\subseteq(R_Z/sR_Z(1))^2\subseteq R_Z/sR_Z(2). \]

记 $W_L=Bt$。则 $t\in W_L$，所以 $1\in B$，且 $B\subset A$ 是余维 1 的 $\C$-向量子空间。上述条件改写为 $Bt^2\subseteq(Bt)^2\subseteq At^2$，即 $Bt^2=B\cdot Bt^2$，从而 $B^2=B$。故 $B$ 是 $A$ 的 $\C$-子代数，而引理 4.53(2)–(3) 意味着：
\begin{enumerate}[label=(\arabic*),leftmargin=2.4em,itemsep=1pt,topsep=2pt]
\item 若 $r\ge 2$，则所有射影 $B\to A\to A/A_j$ 都是满射。
\item 若 $m\neq 0$，则射影 $B\to A\to A/\operatorname{socle}(m)$ 是满射。
\end{enumerate}

由此导出矛盾。首先证明 $B$ 是局部环。否则 $B$ 中存在非零幂等元 $e_1$、$e_2$，满足 $e_1e_2=0$ 且 $e_1+e_2=1$。于是 $B=Be_1\oplus Be_2$、$A=Ae_1\oplus Ae_2$。由 $\operatorname{codim}_AB=1$ 得 $Be_1=Ae_1$ 或 $Be_2=Ae_2$，与上述性质 (1) 矛盾。故 $B$ 是局部环。

其次，由于 $B$ 是满足 $\dim_{\C}B=k-1\ge 2$ 的局部环，注意 $B$ 与 $A$ 都不约化。由性质 (2)，$\Spec A\to\Spec B$ 作为集合的像为单射。特别地，$r=1$。于是可设 $A\cong\C[x]/(x^k)$。再次使用性质 (2)，由 $k-1\ge 2$ 知存在 $\xi=x+cx^{k-1}\in B$（$c\in\C$）。于是 $x=\xi-c\xi^{k-1}\in B$，从而 $A=B$，矛盾。

以上证明了 $R_Z$ 由 $R_Z(1)$ 生成。特别地，当 $n\ge 3$ 时，
% [原文 p.178]
还已知 $R_Z/(s,t)(n)=(R_Z/sR_Z(n))/t(R_Z/sR_Z(n-1))=0$。因此，$R_Z/(s,t)$ 由 $R_Z/(s,t)(1)$ 生成，且其定义理想包含所有次数 $\ge 3$ 的齐次元。这意味着 $R_Z/(s,t)$ 的定义理想 $I$ 由 $I(2)$ 与 $I(3)$ 生成。于是 (2) 由引理 4.55 得到。
\end{proof}

\begin{lemmap}{4.55}
设 $R=\bigoplus_{i\ge 0}R(i)$ 是分次 $\C$-代数，$R(0)=\C$，$s\in R(1)$ 是非零因子。设 $x_1,\dots,x_n\in R$ 是齐次元，给出分次 $\C$-代数 $R/sR$ 的极小生成元组。设 $\phi\colon\C[X_1,\dots,X_n]\to R/sR$ 是由 $X_i\mapsto x_i\bmod sR$ 确定的 $\C$-代数同态。则下列成立。
\begin{enumerate}[label=(\arabic*),leftmargin=2.4em,itemsep=1pt,topsep=2pt]
\item $x_1,\dots,x_n,s$ 是 $R$ 的极小生成元组。特别地，$R$ 是有限生成的。
\item 设 $\psi\colon\C[X_1,\dots,X_n,Y]\to R$ 是由 $X_i\mapsto x_i$、$Y\mapsto s$ 确定的 $\C$-代数同态，则 $(\ker\psi)/Y(\ker\psi)\cong\ker\phi$ 成立。特别地，$\ker\phi$ 的任意极小齐次生成元组可以延拓为 $\ker\psi$ 的极小齐次生成元组。
\end{enumerate}
\end{lemmap}

\begin{proof}
(1) 在各次数上应用中山引理即可得到。(2) 将蛇引理（snake lemma）应用于下面的交换图即可得到。
\begin{tikzcd}
0 \arrow[r] & \C[X,Y] \arrow[r,"Y"] \arrow[d,"\psi"] & \C[X,Y] \arrow[r] \arrow[d,"\psi"] & \C[X] \arrow[r] \arrow[d,"\phi"] & 0 \\
0 \arrow[r] & R \arrow[r,"s"] & R \arrow[r] & R/sR \arrow[r] & 0.
\end{tikzcd}
\end{proof}

\begin{definitionp}{4.56}
首先固定 $w_1,\dots,w_n\in\Z_{>0}$。记 $R:=K[x_1,\dots,x_n]$。对 $R$ 中的单项式定义 $w(\prod x_i^{m_i})=\sum m_iw_i$。更一般地，对 $f\in R$，先把 $f$ 写成单项式之和 $f=\sum a_MM$，再定义 $w(f)=\min_{a_M\neq 0}w(M)$。（约定 $w(0)=\infty$。）由此得到理想 $m^w(n)=\{f\in R\mid w(f)\ge n\}$。$\mathbb{A}^n$ 的权（weight）$w$ 加权的爆破定义为
\[ B_0^w\mathbb{A}^n:=\Proj_R\bigoplus_{n\ge 0}m^w(n). \]
不明确写出权 $w$ 时，有时也简称其为加权爆破（weighted blow up）。
\end{definitionp}
% [Tutor Note: 定义 4.56 尚有一段续文（“对任意 $X\subset\mathbb{A}^n$，$B_0^wX$ 定义为 $X$ 在 $B_0^w\mathbb{A}^n$ 中的双有理变换。□”）起于 p.179，按页界归下块承续，故此处环境先闭合。]
